\documentclass[12pt,a4paper]{article}

%=========================================================
%                   ENCODAGE ET POLICES
%=========================================================
\usepackage[utf8]{inputenc}
\usepackage[T1]{fontenc}
\usepackage[french]{babel}
\usepackage{lmodern}

%=========================================================
%                     MATHÉMATIQUES
%=========================================================
\usepackage{amsmath,amssymb,amsfonts,mathrsfs}
\usepackage{tkz-tab}
\everymath{\displaystyle}

%=========================================================
%                     MISE EN PAGE
%=========================================================
\usepackage[left=1.5cm,right=1.5cm,top=1.5cm,bottom=1.8cm]{geometry}
\usepackage{multicol}
\usepackage{float}
\usepackage{array,multirow}
\usepackage{enumitem}
\usepackage{color, soul}

%=========================================================
%                GRAPHISMES ET COULEURS
%=========================================================
\usepackage{tikz}
\usetikzlibrary{arrows.meta,calc}
\usepackage{pgfplots}
\pgfplotsset{compat=1.15}

\usepackage[most]{tcolorbox}
\tcbuselibrary{skins,hooks}

%=========================================================
%                  LIENS ET ARRIÈRE-PLAN
%=========================================================
\usepackage[colorlinks=true,linkcolor=blue,urlcolor=blue,citecolor=blue]{hyperref}
\usepackage{eso-pic}

%=========================================================
%                   COULEURS ET STYLES
%=========================================================
\definecolor{myred}{RGB}{180, 40, 40}
\definecolor{myblue}{RGB}{20, 50, 120}
\definecolor{mygreen}{RGB}{30, 120, 40}

%=========================================================
%                       FILIGRANE
%=========================================================
\AddToShipoutPicture{
    \AtTextCenter{
        \makebox[0pt]{
            \rotatebox{45}{
                \textcolor[gray]{0.92}{
                    \fontsize{5cm}{5cm}\selectfont PGB
                }
            }
        }
    }
}

\begin{document}

% En-tête
\begin{center}
    \Large\textbf{\underline{\textcolor{myred}{Racine Carrée}}}\\[0.2cm]
    \normalsize\textbf{Prof : M. BA} \hfill \textbf{Classe : Terminale L}\\[0.1cm]
    \textbf{Durée : 6 heures} \hfill \textbf{Année scolaire : 2025 -- 2026}
\end{center}
\hrule
\vspace{0.3cm}

%=========================================================
%    I. La racine carrée d'un nombre réel positif
%=========================================================
\section*{\color{myred}I. La racine carrée d'un nombre réel positif}

\subsection*{\color{myred} Activité 1}
1. Trouver un nombre réel positif dont le carré est égal à $9$.\\
2. Trouver un nombre réel positif dont le carré est égal à $25$.\\
3. Existe-t-il un nombre réel dont le carré est égal à $-16$ ? Justifier.

\subsection*{\color{myred} Définition}
Soit $a$ un nombre réel positif ($a \geqslant 0$).\\
On appelle \textbf{racine carrée} de $a$, et on note $\sqrt{a}$, l'unique nombre réel **positif** dont le carré est égal à $a$.
\[
\sqrt{a} \geqslant 0 \quad \text{et} \quad \left(\sqrt{a}\right)^2 = a
\]

\subsection*{\color{myred} Exemple 1}
\begin{itemize}[label=$\bullet$]
    \item $\sqrt{9} = 3$ car $3 \geqslant 0$ et $3^2 = 9$.
    \item $\sqrt{0} = 0$ et $\sqrt{1} = 1$.
    \item $\sqrt{0{,}25} = 0{,}5$ car $0{,}5^2 = 0{,}25$.
\end{itemize}

\subsection*{\color{myred} Propriété}
Soient $a$ et $b$ deux réels positifs.
\begin{enumerate}
    \item Pour tout réel $a \geqslant 0$, $\sqrt{a^2} = a$.
    \item Pour tout réel $a$, $\sqrt{a^2} = |a|$.
    \item $a = b \iff \sqrt{a} = \sqrt{b}$.
\end{enumerate}

\subsection*{\color{myred} Exemple 2}
\begin{itemize}[label=$\bullet$]
    \item $\sqrt{7^2} = 7$.
    \item $\sqrt{(-5)^2} = |-5| = 5$.
\end{itemize}

\subsection*{\color{myred} Remarque}
La racine carrée d'un nombre négatif **n'existe pas** dans les nombres réels $\mathbb{R}$. L'expression $\sqrt{a}$ n'a de sens que si $a \geqslant 0$.

\subsection*{\color{myred} Application 1}
Calculer les nombres suivants : 
\[
A = \sqrt{49}, \quad B = \sqrt{(-11)^2}, \quad C = \left(\sqrt{13}\right)^2
\]

\subsection*{\color{myred} Correction }
\begin{itemize}[label=$\bullet$]
    \item $A = \sqrt{49} = \sqrt{7^2} = 7$
    \item $B = \sqrt{(-11)^2} = |-11| = 11$
    \item $C = \left(\sqrt{13}\right)^2 = 13$
\end{itemize}

\subsection*{\color{myred} Application 2 }
Déterminer la valeur exacte du réel $D = \sqrt{(\sqrt{3}-2)^2}$.

\subsection*{\color{myred} Solution : }
On a $D = |\sqrt{3}-2|$.\\
Or $1,732 \approx \sqrt{3} < 2$, donc $\sqrt{3} - 2 < 0$.\\
Puisque la valeur absolue d'un nombre négatif est son opposé :
\[
D = -(\sqrt{3}-2) = 2 - \sqrt{3}
\]

%=========================================================
%    II. Racines Carrées Et Opérations
%=========================================================
\section*{\color{myred}II. Racines Carrées Et Opérations}

\subsection*{\color{myred} Activité 2 }
1. Calculer $\sqrt{16 \times 9}$ et $\sqrt{16} \times \sqrt{9}$. Que constate-t-on ?\\
2. Calculer $\sqrt{\frac{100}{25}}$ et $\frac{\sqrt{100}}{\sqrt{25}}$. Que constate-t-on ?\\
3. Calculer $\sqrt{9 + 16}$ et $\sqrt{9} + \sqrt{16}$. Comparer les deux résultats.

\subsection*{\color{myred} 1. Produit de deux racines carrées }

\subsection*{\color{myred} Propriété }
Pour tous réels positifs $a$ et $b$ ($a \geqslant 0$ et $b \geqslant 0$) :
\[
\sqrt{a \times b} = \sqrt{a} \times \sqrt{b}
\]

\subsection*{\color{myred} Exemple }
\begin{itemize}[label=$\bullet$]
    \item $\sqrt{4 \times 25} = \sqrt{4} \times \sqrt{25} = 2 \times 5 = 10$.
    \item Écrire $\sqrt{72}$ sous la forme $a\sqrt{b}$ où $a$ et $b$ sont des entiers avec $b$ le plus petit possible :
    \[
    \sqrt{72} = \sqrt{36 \times 2} = \sqrt{36} \times \sqrt{2} = 6\sqrt{2}
    \]
\end{itemize}

\subsection*{\color{myred} 2. Quotient de deux racines carrées }

\subsection*{\color{myred} Propriété}
Pour tout réel $a \geqslant 0$ et tout réel $b > 0$ :
\[
\sqrt{\frac{a}{b}} = \frac{\sqrt{a}}{\sqrt{b}}
\]

\subsection*{\color{myred} Exemple }
\[
\sqrt{\frac{49}{81}} = \frac{\sqrt{49}}{\sqrt{81}} = \frac{7}{9}
\]

\subsection*{\color{myred} Application 1}
Simplifier les expressions suivantes :
\[
E = \sqrt{50} - 3\sqrt{18} + 2\sqrt{8} \quad \text{et} \quad F = \sqrt{27} \times \sqrt{3}
\]

\subsection*{\color{myred} Correction }
\begin{itemize}[label=$\bullet$]
    \item $E = \sqrt{25 \times 2} - 3\sqrt{9 \times 2} + 2\sqrt{4 \times 2} = 5\sqrt{2} - 3(3\sqrt{2}) + 2(2\sqrt{2}) = (5 - 9 + 4)\sqrt{2} = 0\sqrt{2} = 0$
    \item $F = \sqrt{27 \times 3} = \sqrt{81} = 9$
\end{itemize}

%=========================================================
%    III. Rendre Rationnel Le Dénominateur D'un Quotient
%=========================================================

\section*{\color{myred}III. Rendre Rationnel Le Dénominateur D'un Quotient}

\subsection*{\color{myred} Activité 3 }
Soit le nombre $X = \frac{6}{\sqrt{3}}$.\\
En multipliant le numérateur et le dénominateur par $\sqrt{3}$, écrire $X$ sans radical au dénominateur.

\subsection*{\color{myred} 1. Premier cas : le dénominateur est de la forme $\sqrt{a}$ }

\subsection*{\color{myred} Propriété }
Pour transformer un quotient de la forme $\frac{A}{\sqrt{a}}$ ($a > 0$) en un quotient sans radical au dénominateur, on multiplie le numérateur et le dénominateur par $\sqrt{a}$ :
\[
\frac{A}{\sqrt{a}} = \frac{A \times \sqrt{a}}{\sqrt{a} \times \sqrt{a}} = \frac{A\sqrt{a}}{a}
\]

\subsection*{\color{myred} Exemple }
\[
\frac{5}{\sqrt{2}} = \frac{5 \times \sqrt{2}}{\sqrt{2} \times \sqrt{2}} = \frac{5\sqrt{2}}{2}
\]

\subsection*{\color{myred} 2. Deuxième cas : le dénominateur contient une somme }

\subsection*{\color{myred} Définition}
Les expressions $(a + \sqrt{b})$ et $(a - \sqrt{b})$ sont dites \textbf{expressions conjuguées} l'une de l'autre.\\
Leur produit supprime le radical grâce à l'identité remarquable $(x-y)(x+y) = x^2 - y^2$ :
\[
(a + \sqrt{b})(a - \sqrt{b}) = a^2 - (\sqrt{b})^2 = a^2 - b
\]

\subsection*{\color{myred} Exemple 1}
L'expression conjuguée de $2 + \sqrt{3}$ est $2 - \sqrt{3}$.\\
Leur produit vaut : $(2 + \sqrt{3})(2 - \sqrt{3}) = 2^2 - (\sqrt{3})^2 = 4 - 3 = 1$.

\subsection*{\color{myred} Remarque }
De même, l'expression conjuguée de $\sqrt{a} + \sqrt{b}$ est $\sqrt{a} - \sqrt{b}$.

\subsection*{\color{myred} Exemple 2}
Rendre rationnel le dénominateur de $G = \frac{2}{3 + \sqrt{5}}$ :
\[
G = \frac{2(3 - \sqrt{5})}{(3 + \sqrt{5})(3 - \sqrt{5})} = \frac{2(3 - \sqrt{5})}{3^2 - (\sqrt{5})^2} = \frac{2(3 - \sqrt{5})}{9 - 5} = \frac{2(3 - \sqrt{5})}{4} = \frac{3 - \sqrt{5}}{2}
\]

\subsection*{\color{myred} Exemple 3}
Rendre rationnel le dénominateur de $H = \frac{1}{\sqrt{7} - \sqrt{2}}$ :
\[
H = \frac{1(\sqrt{7} + \sqrt{2})}{(\sqrt{7} - \sqrt{2})(\sqrt{7} + \sqrt{2})} = \frac{\sqrt{7} + \sqrt{2}}{7 - 2} = \frac{\sqrt{7} + \sqrt{2}}{5}
\]

\subsection*{\color{myred} Application }
Écrire les nombres suivants sans radical au dénominateur :
\[
K = \frac{3}{\sqrt{6}} \quad \text{et} \quad L = \frac{4}{\sqrt{5} - 1}
\]

\subsection*{\color{myred} Correction }
\begin{itemize}[label=$\bullet$]
    \item $K = \frac{3\sqrt{6}}{\sqrt{6} \times \sqrt{6}} = \frac{3\sqrt{6}}{6} = \frac{\sqrt{6}}{2}$
    \item $L = \frac{4(\sqrt{5} + 1)}{(\sqrt{5} - 1)(\sqrt{5} + 1)} = \frac{4(\sqrt{5} + 1)}{5 - 1} = \frac{4(\sqrt{5} + 1)}{4} = \sqrt{5} + 1$
\end{itemize}

%=========================================================
%    IV. Résolution de l'équation $x^{2}=a$
%=========================================================

\section*{\color{myred}IV. Résolution de l'équation $x^{2}=a$}

\subsection*{\color{myred} Activité 4 }
Résoudre dans $\mathbb{R}$ les équations suivantes :
1. $x^2 = 16$\\
2. $x^2 = 0$\\
3. $x^2 = -9$

\subsection*{\color{myred} Propriété et Synthèse}
Soit $a$ un nombre réel. Pour résoudre l'équation $x^2 = a$ dans $\mathbb{R}$ :
\begin{itemize}[label=$\bullet$]
    \item Si $a > 0$, l'équation admet **deux solutions réelles distinctes** : $S = \{-\sqrt{a} ; \sqrt{a}\}$.
    \item Si $a = 0$, l'équation admet **une unique solution** : $S = \{0\}$.
    \item Si $a < 0$, l'équation n'admet **aucune solution réelle** : $S = \emptyset$.
\end{itemize}

\subsection*{\color{myred} Exemples}
\begin{enumerate}
    \item $x^2 = 25 \implies x = \sqrt{25}$ ou $x = -\sqrt{25}$, soit $S = \{-5 ; 5\}$.
    \item $x^2 = 7 \implies x = \sqrt{7}$ ou $x = -\sqrt{7}$, soit $S = \{-\sqrt{7} ; \sqrt{7}\}$.
    \item $x^2 = -4 \implies S = \emptyset$ (un carré ne peut pas être négatif dans $\mathbb{R}$).
\end{enumerate}

\end{document}